\chapter{Rings and Modules of Fractions}

\section*{MULTIPLICATIVELY CLOSED SETS}

\begin{proposition}{3.3}
	The operation $S^{-1}$ is exact.
\end{proposition}

\begin{proposition}{3.5}
	Let $M$ be an $A$-module. Then there exists a unique $S^{-1}A$ module isomorphism $f:S^{-1}A\otimes_AM\rightarrow S^{-1}M$.
\end{proposition}

\section*{LOCAL PROPERTIES}

\begin{proposition*}
	Let $P$ be a property of a ring $A$ (or of an $A$-module $M$). Then the following are equivalent:
	\begin{enumerate}[label=\textup{(\roman*)}]
		\item $A$ (or $M$) has $P$ $\Leftrightarrow$ $A_\mathfrak{p}$ (or $M_\mathfrak{p}$) has $P$ for each prime ideal $\mathfrak{p}$ of $A$;
		\item $A$ (or $M$) has $P$ $\Leftrightarrow$ $A_\mathfrak{m}$ (or $M_\mathfrak{m}$) has $P$ for each maximal ideal $\mathfrak{m}$ of $A$.
	\end{enumerate}
	The property $P$ is said to be a \textcolor{orange}{local property}.
	\begin{proof}
		$\left(A_\mathfrak{p}\right)_{\mathfrak{m}A_\mathfrak{p}}\cong A_\mathfrak{m}$
	\end{proof}
\end{proposition*}

\begin{proposition}{3.8}
	Let $M$, $N$ be $A$-modules. Then the following are equivalent:
	\begin{enumerate}[label=\textup{(\roman*)}]
		\item $M=0$;
		\item $M_\mathfrak{p}=0$ for all prime ideals $\mathfrak{p}$ of $A$;
		\item $M_\mathfrak{m}=0$ for all maximal ideals $\mathfrak{m}$ of $A$;
	\end{enumerate}
\end{proposition}

\begin{proposition}{3.10}
	Let $M$ be $A$-modules. Then the following are equivalent:
	\begin{enumerate}[label=\textup{(\roman*)}]
		\item $M$ is a flat $A$-module;
		\item $M_\mathfrak{p}$ is a flat $A_\mathfrak{p}$-module for each prime ideal $\mathfrak{p}$;
		\item $M_\mathfrak{m}$ is a flat $A_\mathfrak{m}$-module for each maximal ideal $\mathfrak{m}$;
	\end{enumerate}
\end{proposition}

\section*{EXTENSIONS AND CONTRACTIONS OF FRACTIONS}

If $\mathfrak{a}$ is an ideal in $A$, its extension $\mathfrak{a}^e$ in $S^{-1}A$ is $S^{-1}\mathfrak{a}$. Denote the contraction of $S^{-1}\mathfrak{a}$ in $A$ to be $S(\mathfrak{a})$.

\begin{proposition}{3.11}
	\begin{enumerate}[label=\textup{(\roman*)}]
		\item Every ideal in $S^{-1}A$ is an extended ideal.
		\item If $\mathfrak{a}$ is an ideal in $A$, then $\mathcolor{orange}{S(\mathfrak{a})}=\mathfrak{a}^{ec}=\bigcup_{s\in S}(\mathfrak{a}:s)$.
		\item $\mathfrak{a}$ is a contracted ideal $\Leftrightarrow$ no element of $S$ is a zero-divisor in $A/\mathfrak{a}$.
		\item Them prime ideal of $S^{-1}A$ are in one-to-one correspondence ($\mathfrak{p}\leftrightarrow S^{-1}\mathfrak{p}$) with the prime ideal of $A$ which don't meet $S$.
		\item The operation $S^{-1}$ commutes with formation of finite sums, products, intersections and radicals.
	\end{enumerate}
\end{proposition}

\section*{DEFINITIONS AND PROPOSITIONS FROM EXERCISES}

\begin{definition}
	A multiplicatively closed subset $S$ of a ring $A$ is said to be \textcolor{orange}{saturated} if 
	\begin{equation*}
		xy\in S\Leftrightarrow x\in S\text{ and }y\in S.
	\end{equation*}
\end{definition}

\begin{definition}
	Let $M$ be an $A$-module. An element $x\in M$ is a \textcolor{orange}{torsion} element of $M$ if there exist a non-zero-divisor $a\in A$ such that $ax=0$. All torsion elements of $M$ form a submodule of $M$, which is called the \textcolor{orange}{torsion submodule} of $M$ and denote by $T(M)$. If $T(M)=0$, the module $M$ is said to be \textcolor{orange}{torsion free}.
\end{definition}

\begin{definition}
	Let $M$ be an $A$-module. The \textcolor{orange}{support} of $M$ is defined to be the set $\operatorname{Supp}(M)$ of prime ideals $\mathfrak{p}$ such that $M_\mathfrak{p}\neq0$.
\end{definition}

\begin{definition}
	$\operatorname{Spec}\left(A_\mathfrak{p}/\mathfrak{p}A_\mathfrak{p}\otimes_AB\right)$ is called the \textcolor{orange}{fiber} of $f^*:\operatorname{Spec}(B)\rightarrow\operatorname{Spec}(A)$ over $\mathfrak{p}$.
\end{definition}
